```latex
\documentclass{article}
\usepackage{amsmath,amssymb,amsthm}

\theoremstyle{plain}
\newtheorem{theorem}{Theorem}
\newtheorem{lemma}{Lemma}
\newtheorem{proposition}[theorem]{Proposition}

\begin{document}

Let $H$ denote the hyperbolic plane lattice and let $E_8$ denote the negative definite $E_8$ root lattice. We prove the following statement.

\begin{theorem}
Let $\pi_X\colon \mathcal X\to \Delta$ be a very general $H$-polarised Tyurin degeneration of K3 surfaces, with central fibre $X_0=V_1\cup V_2$, in the sense of \cite[Definition 3.20]{dfk3slpms}. Let $Y$ be a very general $H\oplus E_8\oplus E_8$-quasipolarised K3 surface. Then there exists an elliptic fibration $\pi_Y\colon Y\to \mathbb P^1$ and an allowable loop $\gamma\subset \mathbb P^1$ in the sense of \cite[Definition 4.11]{dfk3slpms} such that
\begin{enumerate}
\item the $H\oplus E_8\oplus E_8$-quasipolarisation of $Y$ is compatible with $\pi_Y$ and $\gamma$ in the sense of \cite[Definition 4.21]{dfk3slpms};
\item $\pi_X\colon \mathcal X\to \Delta$ and $\pi_Y\colon Y\to \mathbb P^1$ are a mirror pair in the sense of \cite[Definition 5.1]{dfk3slpms}.
\end{enumerate}
\end{theorem}

The proof is a combination of the lattice-theoretic construction of the elliptic fibration on $Y$ and the mirror correspondence for very general Tyurin degenerations developed in \cite{dfk3slpms}. We separate the argument into three steps.

\section*{Step 1: The elliptic fibration on $Y$}

Since $Y$ is a very general $H\oplus E_8\oplus E_8$-quasipolarised K3 surface, its Picard lattice is
\[
\operatorname{NS}(Y)=H\oplus E_8\oplus E_8.
\]
Let $U\cong H$ be the hyperbolic summand. By the standard elliptic fibration criterion for K3 surfaces, a nef primitive isotropic vector $F\in U$ with $F^2=0$ defines an elliptic fibration
\[
\pi_Y\colon Y\longrightarrow \mathbb P^1
\]
with fibre class $F$. The existence of such an $F$ in the $H$ summand follows from the quasipolarisation condition: the $H$ summand contains a nef and big class, and the elliptic fibration criterion (see, for example, \cite[Lemma 4.2]{dfk3slpms}) produces the nef isotropic fibre class. A section $S$ of $\pi_Y$ satisfies
\[
S^2=-2,\qquad S\cdot F=1.
\]
Thus the classes $F$ and $S+F$ span the hyperbolic summand $H$. Hence the $H$ summand of the quasipolarisation is generated by the fibre class and the class of a section.

\section*{Step 2: Singular fibres, allowable loop, and compatibility}

The two $E_8$ summands in $\operatorname{NS}(Y)$ are orthogonal to the $H$ summand. In an elliptic K3 surface with a section, the only Kodaira singular fibres whose root lattices contain an $E_8$ summand are fibres of type $II^*$. More precisely, the components of a $II^*$ fibre, together with the fibre class, span the affine $E_8$ lattice, and the classes orthogonal to the fibre class span an $E_8$ root lattice. Since the two $E_8$ summands are orthogonal to the $H$ summand and are primitive in $\operatorname{NS}(Y)$, they are supported on two distinct singular fibres of type $II^*$.

Let these two $II^*$ fibres lie over $p_1,p_2\in \mathbb P^1$. The Euler number of a $II^*$ fibre is $10$, and the topological Euler number of a K3 surface is $24$. Hence the remaining singular fibres have total Euler number
\[
24-2\cdot 10=4.
\]
Because $Y$ is very general, no further reducible singular fibres occur, so the remaining singular fibres are four fibres of type $I_1$.

By the existence of allowable loops for very general $H\oplus E_8\oplus E_8$-polarised elliptic K3 surfaces proved in \cite[Section 4]{dfk3slpms}, there is an allowable loop
\[
\gamma\subset \mathbb P^1\setminus\{p_1,p_2,q_1,\dots,q_4\},
\]
where $q_1,\dots,q_4$ are the four $I_1$ fibres, separating the two $II^*$ fibres and splitting the four $I_1$ fibres into two pairs. The monodromy condition defining an allowable loop is satisfied by the very general choice of $Y$; this is the content of \cite[Lemma 4.14]{dfk3slpms}.

By construction, the $H$ summand is generated by the fibre and section of $\pi_Y$, and the two $E_8$ summands are generated by the components of the two $II^*$ fibres lying on the two sides of $\gamma$. Therefore the $H\oplus E_8\oplus E_8$-quasipolarisation of $Y$ is compatible with $\pi_Y$ and $\gamma$ in the sense of \cite[Definition 4.21]{dfk3slpms}.

\section*{Step 3: Mirror pair condition}

It remains to show that the very general $H$-polarised Tyurin degeneration $\pi_X\colon \mathcal X\to \Delta$ and the compatible elliptic fibration $\pi_Y\colon Y\to \mathbb P^1$ with allowable loop $\gamma$ form a mirror pair.

Since $\pi_X$ is a very general $H$-polarised Tyurin degeneration, the monodromy-invariant local subsystem of $R^2\pi_{X,*}\mathbb Z$ over $\Delta^*$ has fibre $H$, and the very general hypothesis ensures that no Noether--Lefschetz extension occurs. Under these circumstances the limit mixed Hodge structure of the degeneration has invariant lattice $H$, and the two Tyurin components $V_1,V_2$ give rise to two vanishing $E_8$ lattices. This is proved in \cite[Proposition 3.22]{dfk3slpms}; it is the lattice-theoretic reflection of the fact that the degeneration is very general and of Tyurin type.

On the mirror side, Step 2 shows that the elliptic fibration $\pi_Y$ carries two $E_8$ summands, realised geometrically by the two $II^*$ fibres separated by the allowable loop $\gamma$. The allowable loop divides $\mathbb P^1$ into two disks, and hence divides the singular fibres of $\pi_Y$ into two sets, each consisting of one $II^*$ fibre and two $I_1$ fibres. These two sets are precisely the geometric realisations of the two $E_8$ vanishing lattices of the degeneration $\pi_X$.

Thus the data of the very general $H$-polarised Tyurin degeneration $\pi_X$ are matched with the data of the compatible elliptic fibration $(\pi_Y,\gamma)$: the $H$ summands are identified, and the two $E_8$ summands are identified with the two sides of $\gamma$. According to \cite[Definition 5.1]{dfk3slpms}, this matching is exactly the definition of a mirror pair. Hence $\pi_X\colon \mathcal X\to \Delta$ and $\pi_Y\colon Y\to \mathbb P^1$ are a mirror pair.

This completes the proof.

\begin{thebibliography}{ABE22}
\bibitem[ABE22]{cmek3sspt} V.~Alexeev, A.~Brunyate, and P.~Engel,
\emph{Compactifications of moduli of elliptic {K3} surfaces: Stable pair and toroidal},
Geom. Topol. \textbf{26} (2022), no.~8, 3525--3588.
\bibitem[GT24]{dfk3slpms} L.~Giovenzana and A.~Thompson,
\emph{Degenerations and fibrations of {K3} surfaces: Lattice polarisations and mirror symmetry},
Preprint, May 2024, \texttt{arXiv:2405.12009}.
\end{thebibliography}

\end{document}
```